% Bu dosya, Bolum 9'u zenginlestirir (Yeni chapter acmaz)
\section{IR Playbook'lar ve Triyaj Kontrol Listeleri}

\subsection{Playbook Özeti}
\begin{longtable}{|p{3.3cm}|p{5.5cm}|p{6cm}|}
\hline
\rowcolor{tableheadcolor}
\textbf{Senaryo} & \textbf{Adımlar (SANS Aşamaları)} & \textbf{Notlar} \\
\hline
Fidye Yazılımı & Tespit→Kısa Sınırlama→Adli Kopya→Ağ Segmentasyonu→Yedekten Dönüş & Ödeme kararı için hukuk ve üst yönetim bilgilendirilmeli \\
\hline
Oltalama (Kimlik İhlali) & Tespit→Hesap Dondurma→Parola Sıfırlama→OAuth Token İptali→Kayıt İncelemesi & Kullanıcıya bilgilendirme; MFA zorunluluğu \\
\hline
Web Shell & Tespit→WAF Kuralı→İlgili Host İzolasyonu→Log/Adli Toplama→Yama ve Sertleştirme & Gerekirse dış IR desteği alın \\
\hline
Bulut İçi Anahtar Sızıntısı & Tespit→Erişim Anahtarı İptali→IAM Politika İncelemesi→Envanter Taraması & CI/CD sırları ve repo geçmişi taranmalı \\
\hline
\end{longtable}

\subsection{Forensik Triyaj Kontrol Listesi}
\begin{itemize}
  \item \textbf{Zaman Çizgisi:} Sistem saatini kaydet, olay saat penceresi
  \item \textbf{Uç Nokta:} Çalışan işlemler, ağ bağlantıları, kalıcılık noktaları (Run Keys, Services)
  \item \textbf{Ağ:} PCAP, proxy/DNS kayıtları, egress kontrol
  \item \textbf{Bulut:} Yönetim aksiyonları, IAM değişiklikleri, token ve rol atamaları
\end{itemize}

\Not{Playbook'lar yaşayan dokümanlardır. Post-mortem sonuçlarına göre düzenli olarak güncellenmelidir.}

\Ipucu{Tüm adımlar için case/etiket standardı kullanın: olay türü, etkilenen varlık, ATT\&CK tekniği, CVE/CWE (varsa).}

